\documentclass[openany]{book}

\usepackage{multicol}
\usepackage[dvipsnames]{xcolor}
\usepackage{lettrine}
\usepackage{geometry}
\geometry{
% Try these
	paperwidth=5.0in,
	paperheight=7.5in,% Could be taller, 7.75
	twoside,
	%
	% --- HORIZONTAL MARGINS ---
	% KDP minimum for a ~500 page book is 0.625in.
	% We use 0.65in to protect the text from disappearing into the deep glue spine.
	% inner=0.65in,
	% inner=0.75in,
	inner=0.765in,
	% inner=0.85in,
	% KDP minimum is 0.25in. We use 0.3in to give the text a beautiful,
	% balanced frame on the wider 5.0" page.
	outer=0.30in,
	%
	% --- VERTICAL MARGINS ---
	% Since you removed the tall header, 0.35in is tight, clean, and safe.
	% top=0.35in,
	top=0.28in,
	% Gives your footer page number breathing room above the 0.25" KDP cutting floor.
	% bottom=0.35in,
	% bottom=0.28in,
	% bottom=0.26in,
	% bottom=0.25in,
	bottom=0.2in,
	%
	% --- MEASUREMENT OVERRIDES ---
	% Tells geometry to compute margins from the top of the actual header line
	includehead,
	nomarginpar,
    % Explicitly tightens the gap between text baseline and the page number
    footskip=14pt,
    % footskip=18.5pt,
includehead,           % Keeps the tall header inside the 0.45in top safety zone
includefoot            % Keeps the footer inside the 0.4in bottom safety zone
}

\usepackage{fancyhdr}

\clubpenalty=10000

\newcommand{\rank}{\textit}
\newcommand{\antiphon}{\textit}

\begin{document}

\thispagestyle{fancy}

\begin{center}
\begin{tabular}{l|p{3in}}
1  & Dom. XXII. P. T. \textbf{Festivitas Omnium Sanctorum.} \rank{T. D.} \\
2  & Fer. 2. \\
3  & Fer. 3. \\
4  & Fer. 4. \\
5  & Fer. 5. \\
6  & Fer. 6. \\
7  & Sab. \\
8  & Dom. \\
9  & Fer. 2. \\
10 & Fer. 3. \\
11 & Fer. 4. \\
12 & Fer. 5. \\
13 & Fer. 6. \\
14 & Sab. \\
15 & Dom. \\
16 & Fer. 2. \\
17 & Fer. 3. \\
18 & Fer. 4. \\
19 & Fer. 5. \\
20 & Fer. 6. \\
21 & Sab. \\
22 & Dom. \\
23 & Fer. 2. \\
24 & Fer. 3. \\
25 & Fer. 4. \\
26 & Fer. 5. \\
27 & Fer. 6. \\
28 & Sab. \\
29 & Dom. \\
30 & Fer. 2. \\
\end{tabular}
\end{center}


\end{document}